% =====================================================================
%  AS.410.662 Molecular Modeling for Drug Discovery Lab
%  RESEARCH PROJECT - PREPRINT MANUSCRIPT
%  Rayan Abdelaziz
%  Built on the course Report Template. Formatting is unmodified.
% =====================================================================

\documentclass[sigconf]{acmart}

\usepackage{graphicx}
\usepackage{amsmath}
\usepackage{float}
\usepackage{comment}
\usepackage{textcomp}
\usepackage{makecell}
\usepackage{booktabs}
\usepackage{listings}
\usepackage{xcolor}

\settopmatter{printacmref=false} % removes the citation block below the abstract
\renewcommand\footnotetextcopyrightpermission[1]{} % removes the copyright footnote
\pagestyle{plain}                % removes running headers

% Marks a number still pending the full-ensemble validity sweep.
% Must be empty before submission.
\newcommand{\pending}[1]{\textbf{[PENDING: #1]}}

\begin{document}

\title{Receptor Identity, Not Cognate Pairing, Governs Docking Success
       at the $\alpha$7 Nicotinic Receptor Modulator Site}

\author{Rayan Abdelaziz}
\email{rabdela7@jh.edu}
\orcid{}
\affiliation{%
  \institution{Johns Hopkins University}
  \city{Baltimore}
  \state{Maryland}
  \country{USA}
}

\begin{abstract}
The $\alpha$7 nicotinic acetylcholine receptor is a cholinergic target whose
positive allosteric modulators (PAMs) avoid the desensitisation that has
frustrated orthosteric agonists. Four cryo-EM structures now resolve four
chemically distinct PAMs in the same transmembrane site, which makes it possible
to ask how much a docking protocol depends on which of them is used. Each PAM was
docked into each structure across five seeds in one shared coordinate frame and
one shared box, giving 80 runs per preparation arm and 160 in total. Twenty-six of
80 runs placed a correct pose at rank 1, and every one of those 26 occurred in
receptor 8V80 or 8V88; 8UTB and 8V82 returned none, including for their own
ligands. Success varied far more across receptors (0, 11, 0, 15 of 20) than across
ligands (0, 6, 10, 10 of 20), and the cognate diagonal split two cells at 5/5
against two at 0/5, so cognate docking was not an upper bound. The two failures
were mechanistically distinct: 8V82 sampled a 0.71--0.79~\AA\ pose in every seed
and ranked it second, whereas 8UTB never sampled one. Of the 720 poses, 715 passed
all 22 physical-validity checks and the 5 that failed were low-ranked poses of a
single ligand, so misplacement rather than chemistry drives the result. A second,
independently prepared arm using the recommended p$K_a$-based protonation
reproduced the rank-1 matrix cell for cell, excluding protonation as the
explanation.
\end{abstract}

\keywords{Molecular Docking, Cross-Docking, Structure-Based Drug Design,
Allosteric Modulation, Nicotinic Acetylcholine Receptor, Pose Prediction}

\maketitle
\fancyfoot{}
\fancyhead{}
\thispagestyle{empty}

% =====================================================================
\section{Introduction}
\label{sec:introduction}

The $\alpha$7 nicotinic acetylcholine receptor is a homopentameric Cys-loop
channel widely expressed in the nervous system and on immune cells, and a
long-standing target in cognitive and inflammatory disease. Orthosteric agonists
have repeatedly failed to translate, largely because the receptor desensitises
within milliseconds, so attention has moved to positive allosteric modulators
(PAMs), which act only when acetylcholine is already present. PAMs bind an
intersubunit pocket in the transmembrane domain, so the site is formed by two
adjacent subunits and is sensitive to their relative orientation.

Burke et al.~\cite{burke2024} determined cryo-EM structures of $\alpha$7 bound to
epibatidine together with each of several chemically distinct PAMs, and reported
that the modulators occupy partially overlapping sites while the receptor adopts
different conformational states by modulator class. If one site adopts materially
different shapes depending on which modulator is bound, then any virtual screen
built on one of these structures inherits that structure's pocket. Cross-docking
performance has been characterised for soluble enzyme families, but not for this
site, and not for a set of structures differing only in bound modulator; which of
the four a screen should use is an unexamined decision.

\textbf{Hypothesis.} This study was proposed to test a specific prediction: that
cognate receptor conformations recover their own experimentally observed PAM poses
more accurately than noncognate conformations. That prediction follows directly
from the cross-docking literature on soluble enzymes, and it is the hypothesis the
experiment was designed around. It is not supported by the result, and
Section~\ref{sec:hypothesis} takes up why.

\textbf{Contribution.} We run the full cognate/cross matrix for the four PAM-bound
structures in a single shared frame with seed replication; decompose each failure
into sampling versus scoring using the full ranked pose list rather than the top
pose alone; and test the leading preparation-based explanation by repeating the
entire matrix under a second protonation protocol differing in one variable.

\subsection{Related Work}

Sutherland et al.~\cite{sutherland2007} cross-docked eight protein families and
established the two findings this study is positioned against: average
cross-docking accuracy is substantially worse than self-docking, and accuracy
improves when the receptor is chosen from the complex whose bound ligand most
resembles the docked ligand. They also found no relationship between
crystallographic resolution and cross-docking accuracy. Their targets were soluble
enzymes with one ligand per site; whether the same holds at a lipid-exposed
intersubunit allosteric site in a membrane channel is open. Amaro et
al.~\cite{amaro2018} review ensemble docking, the standard remedy once receptor
dependence is established, which presupposes that conformational variation drives
the variance --- testable here, and measured directly.

Buttenschoen et al.~\cite{buttenschoen2024} introduced PoseBusters after showing
that methods can report confident poses that are not physically valid molecules,
and argued physical-validity checks should accompany any pose claim. We apply
those checks and find they separate cleanly from pose correctness in this system,
which bounds what such filtering can be expected to do. What none of this work
addresses is whether the cognate/cross distinction is the right axis at all ---
the question the matrix below answers in the negative.

% =====================================================================
\section{Methods}

\subsection{Structures and Ligand Set}

Four PAM-bound $\alpha$7 structures from Burke et al.~\cite{burke2024}, retrieved
from the PDB~\cite{berman2000pdb} on 2026-11-11: 8UTB (NS-1738, ligand XG3,
2.30~\AA), 8V80 ((--)-TQS, YLR, 2.34~\AA), 8V82 (PNU-120596, I34, 2.61~\AA) and
8V88 (GAT107, YLI, 2.30~\AA). All four carry the same orthosteric agonist
(epibatidine), so they differ in modulator and conformational state rather than in
construct. $\alpha$7 is homopentameric, so each structure holds up to five copies
of the PAM site; the copy at the A--B interface was selected by measuring each
copy's distance to that interface rather than by hard-coding a chain and residue,
which independently reproduced the copy chosen manually in earlier work on 8V82.

\subsection{Common Frame and Docking Box}

Results are comparable across receptors only if the box is the same box. Chains A
and B of each structure were matched to 8V82 by residue number and superposed on
matched C$\alpha$ atoms with \texttt{gemmi}~\cite{wojdyr2022}, applying the
transform to the whole structure including its ligand, so each receptor and its
own crystallographic pose move together. One box was then defined once in that
frame --- centre $(146.903, 153.799, 133.568)$~\AA, size $20.0 \times 22.0 \times
29.0$~\AA, 12{,}760~\AA$^3$ --- and used for all 160 runs. No per-receptor
recentring was done, since that would reintroduce the receptor-specific bias the
common frame removes.

\subsection{Receptor Preparation: Two Arms}

Preparation was run as two arms differing in exactly one variable.
\textbf{Arm~1} repaired structures with PDBFixer (missing heavy atoms and residues
rebuilt, hydrogens added at pH~7.4) via the OpenMM toolchain~\cite{eastman2017},
then converted to PDBQT with Meeko~\cite{meeko}. \textbf{Arm~2} was identical
through heavy-atom repair, after which protonation was assigned by PDB2PQR with
the AMBER force field~\cite{dolinsky2007} and titration states predicted by PROPKA
at pH~7.4~\cite{olsson2011}, recording each histidine state (HID/HIE/HIP) in the
run manifest. Arm~2 is the recommended route for this toolchain, because hydrogens
placed by PDBFixer are a known source of valence errors when Meeko reads the
receptor; arm~1 is the alternative. Lipids, waters, ions, glycans and epibatidine
were excluded uniformly from all four receptors in both arms, so no
cognate-versus-cross difference can reflect a per-receptor preparation choice.

One preparation detail changed the protocol and is reported. Meeko rejected a
C-terminal residue in 8V80 with an explicit-valence error traceable to the OXT
atom added by PDBFixer --- an instance of the incompatibility above --- and its
own options for discarding problem residues act after the failing stage, so the
fix was to strip terminal OXT atoms before Meeko reads the structure. A guard was
added that refuses to drop any unbuildable residue within 8.0~\AA\ of the box;
arm~2 later tripped it at 8V80 B:478 (1.11~\AA\ from the box) for the identical
reason. The manifest records \texttt{dropped\_unbuildable\_residues: []} for every
receptor in both arms: no binding-site residue was ever silently discarded. The
same OXT rule was applied in both arms, preserving the single-variable difference.

\subsection{Ligand Preparation and Docking}

Ligand bond orders and formal charges were taken from the RCSB chemical component
dictionary definition for each code, never perceived from deposited coordinates,
since geometry-based perception is unreliable for the aromatic and sulfonamide
groups these PAMs carry and a wrong bond order changes the rotatable-bond count
and therefore the search problem. Ligands were built with RDKit~\cite{rdkit} and
converted with Meeko~\cite{meeko}. Docking used AutoDock Vina
1.2.7~\cite{eberhardt2021autodock,trott2010autodock} at exhaustiveness 32, 9
modes, 2 CPU threads, with seeds 2026--2030: $4 \times 4 \times 5 = 80$ runs and
up to 720 poses per arm.

\subsection{Validation, Analysis and Metrics}

A regression check was defined before the matrix and required to pass before any
matrix run started: the generalised pipeline had to reproduce a previously
established single-receptor 8V82 baseline, rank-1 RMSD in $[6.40, 6.46]$~\AA\ and
rank-2 in $[0.70, 0.80]$~\AA, with the decision rule set in advance. This guards
against the characteristic failure of a refactor --- a pipeline that silently
changes behaviour and returns a clean matrix of wrong numbers. Arm~1 returned
6.408 / 0.744~\AA\ and arm~2 6.443 / 0.723~\AA; both passed.

Pose accuracy is symmetry-aware heavy-atom RMSD to the crystallographic pose via
RDKit \texttt{CalcRMS}~\cite{rdkit}, computed in place. Symmetry awareness is
required because these ligands contain topologically equivalent atoms; no prior
alignment is applied because both poses already sit in the common frame, and
aligning would measure shape rather than placement agreement. Two measures are
reported, and the distinction carries the analysis: \textbf{top-1 success}, the
rank-1 pose within 2.0~\AA\ --- what a screen consumes --- and \textbf{any-pose
success}, any pose in the ranked list within 2.0~\AA\ --- whether the search ever
visited the answer. Failing top-1 while passing any-pose is a \emph{scoring}
failure; failing both is a \emph{sampling} failure. A single combined rate
discards exactly that distinction.

Physical validity used PoseBusters~\cite{buttenschoen2024} (22 checks: internal
geometry, stereochemistry, intramolecular and protein--ligand clashes), invoked
without a reference ligand, so it contributes validity checks only and no RMSD
check. Pose diversity used DBSCAN~\cite{ester1996} on the precomputed pairwise
RMSD matrix within each condition (45 poses: 9 modes $\times$ 5 seeds),
$\varepsilon = 2.0$~\AA, minimum cluster size 3 --- clustering on the RMSD matrix
directly keeps the clustering metric identical to the success metric. Native
contacts are heavy-atom proximities within 4.5~\AA\ of the deposited pose, and are
proximities only: no donor--acceptor angle is tested.

% =====================================================================
\section{Results}

\subsection{Success Is Organised by Receptor}

Table~\ref{tab:matrix} gives arm-1 outcomes for all 16 cells.

\begin{table}[h!]
\caption{Arm 1: top-1 / any-pose successes out of 5 seeds. Cognate cells marked
$\dagger$. Every top-1 success in the entire matrix falls in receptor 8V80 or
8V88.}
\label{tab:matrix}
\begin{center}
\begin{tabular}{l|c|c|c|c}
\hline
Ligand & 8UTB & 8V80 & 8V82 & 8V88 \\ \hline
8UTB (NS-1738)    & 0/0$^\dagger$ & 1/5 & 0/0 & 5/5 \\ \hline
8V80 ((--)-TQS)   & 0/0 & 5/5$^\dagger$ & 0/0 & 5/5 \\ \hline
8V82 (PNU-120596) & 0/3 & 0/0 & 0/5$^\dagger$ & 0/0 \\ \hline
8V88 (GAT107)     & 0/0 & 5/5 & 0/0 & 5/5$^\dagger$ \\ \hline
Receptor total    & 0/3 & 11/15 & 0/5 & 15/15 \\ \hline
\end{tabular}
\end{center}
\end{table}

Across 80 runs, 26 placed a correct pose at rank 1 and 38 placed one anywhere.
All 26 top-1 successes occurred in 8V80 or 8V88; 8UTB and 8V82 returned 0/20 each,
for every ligand including their own. Summed over ligands the four receptors gave
0, 11, 0 and 15 of 20; summed over receptors the four ligands gave 0, 6, 10 and 10
of 20. The receptor axis spans the full range and the ligand axis does not.

The cognate diagonal is split: 8V88 succeeds 5/5 (0.313--0.332~\AA) and 8V80 5/5
(0.791--0.828~\AA), while 8V82 is 0/5 at top-1 and 8UTB 0/5 on both measures.
Pooled, the diagonal gives 10/20 cognate against 16/60 cross, with median rank-1
RMSD 3.615 against 7.011~\AA --- numbers that read as a cognate advantage but are
produced by which receptors sit on the diagonal. The sharpest counterexample is
NS-1738: 0/5 in its own 8UTB, with nothing closer than 5.836~\AA, and 5/5 at
0.676--0.689~\AA\ in the noncognate 8V88.

\subsection{Sampling and Scoring Failures Are Distinct}

Cognate 8V82 and cognate 8UTB both report 0/5 at top-1 by different mechanisms. In
8V82 all five seeds place a pose at 0.712--0.786~\AA, always at rank 2, while
rank 1 sits at 6.401--6.458~\AA: the search finds the pose every time and the
scoring function ranks a wrong one above it every time. In 8UTB no seed produces
anything within 2.0~\AA, the best across five seeds being 5.836--5.898~\AA\ at
rank 3. The first is recoverable by rescoring, with the correct pose already in
the output file; the second is not recoverable at all.

\subsection{Wrong Poses Are Chemically Valid}

All 720 poses of arm 1 were checked, 80 at each of ranks 1 through 9.
\textbf{715 passed all 22 checks and 5 did not.} Every failure is a GAT107 pose
(four in receptor 8V88, one in 8V80); every one fails the same single check,
minimum distance to protein, at 3.60--4.12~\AA\ RMSD; and \emph{none is a rank-1
pose} --- they sit at ranks 3, 3, 4, 8 and 9. That is why the full ensemble had to
be checked rather than a subset: an earlier pass over each run's rank-1 and
best-sampled pose covered 130 distinct poses and returned 130 of 130 passing,
which would have supported a stronger claim than the data do, because all five
failures lie at ranks that subset never sampled.

The conclusion for reported poses is unchanged and now better supported. No pose
this protocol would hand to a virtual screen is physically invalid, and the wrong
ones are reasonable molecules in the wrong place. Physical validity therefore does
not discriminate among the poses whose placement is at issue, and validity
filtering cannot serve as a correctness filter for conventional docking output.

Two further measures behave the same way. DBSCAN shows the 8V82 near-native pose
is reproducible rather than a fluctuation: it forms its own cluster, populated in
all five seeds with no points labelled noise, while the incorrectly top-ranked
pose sits in the largest cluster of that condition (14 of 45) --- the scoring
function prefers a densely sampled wrong basin to a reproducible right one. Native
contact recovery averages 92.6\% for near-native poses against 56.4\% for
non-native, but the 8V82 rank-1 pose recovers 73.3\% while being 6.4~\AA\ wrong,
so contact recovery is not a proxy for pose accuracy either.

\subsection{Structural Correlates}
\label{sec:correlates}

8UTB differs from the other three on three independent counts. In the common
frame, its matched C$\alpha$ atoms deviate from 8V82 by 1.455~\AA\ overall and
1.99~\AA\ over the pocket, against 0.379/0.576~\AA\ (8V80) and 0.927/0.610~\AA\
(8V88); superposing the 34-residue pocket locally, 8UTB differs from the others by
1.807, 1.802 and 1.824~\AA\ while all other pairs fall between 0.260 and
0.506~\AA. It is the only structure with no lipid within 6~\AA\ of any PAM copy
(nearest 8.71~\AA), where the other three each place two lipids against every copy
(nearest 2.26--2.42~\AA, second 2.87--3.06~\AA). And it contains no bound
Ca$^{2+}$, while the others hold five copies each. 8V82, by contrast, is not a
conformational outlier; it is the lowest-resolution structure (2.61 versus
2.30--2.34~\AA), and its failure is a scoring failure.

\subsection{Protonation Does Not Explain It}

PROPKA assigns HIS114 as HIE in both failing receptors and HID in both succeeding
ones --- a pattern that arises by chance roughly one time in three at $n = 4$, and
so is testable rather than evidential. Arm~2 is that test. It returned 26/80
top-1, identical to arm~1, and 41/80 any-pose against 38. Compared cell by cell
over all 80 matched runs, \emph{no cell changed its top-1 count}: the rank-1
matrix is identical under both protocols. The three extra any-pose recoveries fall
in two cells, both in receptor 8UTB (cognate 0/5~$\rightarrow$~1/5; 8V82 into 8UTB
3/5~$\rightarrow$~5/5).

An identical matrix invites the objection that the second arm changed nothing, so
the runs were compared numerically rather than by outcome. Of the 80 matched runs,
41 are identical in rank-1 RMSD and score to within $10^{-6}$, and 39 differ. The
39 that differ are \emph{entirely} in the two failing receptors: all 20 runs in
8UTB and 19 of 20 in 8V82. Both succeeding receptors are unchanged, 20 of 20 each,
because the two protocols assigned them equivalent protonation and so produced
equivalent prepared receptors. The negative result is therefore stronger than the
identical matrix alone suggests: changing the protonation protocol altered the
numerical outcome of nearly half the matrix, concentrated precisely on the two
receptors that fail, and did not move a single run across the threshold a virtual
screen would use. Protonation perturbs the search where the search is failing, and
does not rescue it.

% =====================================================================
\section{Discussion}

\subsection*{The hypothesis was wrong, and that is the finding}
\label{sec:hypothesis}

The proposal predicted that cognate conformations would recover their own poses
more accurately than noncognate ones. The matrix does not support that. Two of the
four cognate pairs fail in every seed, one ligand docks better into a foreign
receptor than into its own, and the variable that separates success from failure
is which receptor was used: an otherwise identical protocol returns 15/20 at top-1
in one receptor and 0/20 in another, over the same ligands and the same box. The
prediction was a reasonable one --- it is what the cross-docking literature on
soluble enzyme families would lead you to expect~\cite{sutherland2007} --- and it
does not hold at this site, for these structures.

It is worth being precise about what failed. The hypothesis assumed the cognate
relationship is the informative one: that a pocket shaped around a ligand is
thereby the best pocket for it. Here, two pockets are good for every ligand and
two are good for none, and that property of the receptor swamps the cognate
relationship entirely. The experiment was built to measure a cognate-versus-cross
gap and instead measured a receptor effect the cognate framing cannot express ---
which is why the pooled cognate rate (10/20 against 16/60) looks like a textbook
confirmation while every individual cell contradicts it. Had the matrix been
summarised the way such studies usually are, the wrong conclusion would have been
well supported by the numbers. Cognate docking is normally treated as a protocol's
ceiling; here a pooled rate implies a moderate capability that no individual cell
exhibits, since two cells succeed completely, two fail completely, and the mean
describes neither. For anyone screening against this site, which deposited
structure to use is a larger decision than ligand preparation or search depth.

This contradicts the usual reading of Sutherland et al.~\cite{sutherland2007}:
their similarity-selection result --- pick the receptor whose bound ligand
resembles yours --- presumes the cognate relationship is informative, and here it
is not. That the lowest-resolution structure is also a failing one is likewise not
evidence that resolution predicts accuracy; they found no such relationship, and
with $n = 4$ ours cannot establish one. Theirs is the better-powered finding, and
we read our result as consistent with it.

The sampling/scoring decomposition is worth retaining for the same reason. Cognate
8V82 and cognate 8UTB are identical at top-1 and are different problems with
different remedies --- rescoring for one, and for the other, the conformational
sampling that ensemble docking~\cite{amaro2018} exists to supply. That 8UTB is
simultaneously the conformational outlier, the only lipid-free pocket and the only
Ca$^{2+}$-free structure makes it the natural ensemble-docking test case, though
at $n = 4$ these three differences cannot be disentangled and none is established
as causal. Finally, the validity results bound what physical-validity filtering
can do. Poses 6--10~\AA\ from the answer are valid molecules in plausible
environments, and one recovers nearly three-quarters of native contacts. Such
filtering is an effective guard against the distorted output that motivated
PoseBusters~\cite{buttenschoen2024}, but it does not discriminate among
conventional docking poses and should not be reported as though it does.

\textbf{Limitations.} Two bear directly on the central claim. Three of four
receptors are near-identical in the pocket (0.260--0.506~\AA\ pairwise), so the
conformational arm rests effectively on one structure, and a receptor effect
across four structures of which three are interchangeable is weaker than the count
suggests. And the deposited PAM orientation in 8UTB is the least certain of the
four, modelled into comparatively diffuse density; every 8UTB RMSD is measured
against that reference, so a sampling failure diagnosed there is not fully
separable from disagreement with the reference itself. The 8V82 scoring failure
carries no such caveat. Beyond these: all four structures come from one study, so
the conformational range is that study's; receptors are rigid, which is the
standard Vina setting and precisely the assumption cross-docking stresses; scoring
is Vina's function alone; five seeds bound run-to-run variation but do not support
distributional claims; and the 2.0~\AA\ threshold is conventional but arbitrary,
so RMSD ranges are given throughout. There is no experimental validation.

% =====================================================================
\section{Conclusion}

Across a seed-replicated 4$\times$4 cognate/cross matrix in one shared frame,
docking success at the $\alpha$7 PAM site was determined by which receptor was
used, not by whether a ligand was returned to its own structure: all 26 top-1
successes in 80 runs fell in two of four receptors, and the cognate diagonal split
evenly between complete success and complete failure. The two failing receptors
failed differently --- one never sampling the correct pose, one sampling it in
every seed and ranking it second --- and that distinction is invisible to a top-1
rate. A second matrix under the recommended protonation protocol reproduced the
rank-1 result cell for cell, excluding preparation as the explanation. The
experiment that would most change this conclusion is flexible-receptor or ensemble
docking against 8UTB, which would test directly whether its sampling failure is
the induced-fit problem its conformational, lipid and Ca$^{2+}$ differences
suggest.

% =====================================================================
\section*{Data and Code Availability}

Structures were retrieved from the Protein Data Bank (8UTB, 8V80, 8V82, 8V88;
accessed 2026-11-11) and ligand definitions from the RCSB chemical component
dictionary (XG3, YLR, I34, YLI; same date). All docking, analysis and figure code
is provided as a Jupyter notebook that runs top to bottom on a free Google Colab
CPU runtime, building an isolated Python 3.12 environment with pinned package
versions (AutoDock Vina 1.2.7, Meeko 0.8.0, gemmi 0.7.4, RDKit 2025.9.4, PDBFixer
1.12.0, OpenMM 8.3.1, NumPy 1.26.4, PoseBusters 0.6.5, PDB2PQR 3.8.0, PROPKA
3.5.1), recorded seeds (2026--2030) and a provenance cell that prints the resolved
versions. The notebook reproduces every figure and table from archived per-run
outputs by default and can re-run the docking from a single flag; it reads nothing
from any personal cloud storage. Per-run checkpoints, all 1{,}440 poses, analysis
tables and the run manifest with per-receptor SHA-256 hashes are archived
alongside it.

% =====================================================================
\bibliographystyle{ACM-Reference-Format}
\bibliography{references}

\end{document}
